% Generated draft: requires mathematical and source review.
\documentclass[11pt]{article}
\usepackage[margin=25mm]{geometry}
\usepackage{fontspec}
\setmainfont{Latin Modern Roman}
\usepackage{amsmath,amssymb,mathtools,unicode-math}
\setmathfont{Latin Modern Math}
\usepackage{enumitem,needspace,xurl,hyperref}
\setlength{\parindent}{0pt}
\setlength{\parskip}{5pt}
\setlength{\emergencystretch}{3em}
\newcommand{\N}{\mathbb{N}}
\newcommand{\Z}{\mathbb{Z}}
\newcommand{\Q}{\mathbb{Q}}
\newcommand{\R}{\mathbb{R}}
\newcommand{\C}{\mathbb{C}}
\newcommand{\F}{\mathbb{F}}
\newcommand{\A}{\mathbb{A}}
\newcommand{\PP}{\mathbb{P}}
\newcommand{\E}{\mathbb{E}}
\newcommand{\eps}{\varepsilon}
\newcommand{\dd}{\,\mathrm{d}}
\begin{document}
\section*{UnsolvedMath ID: 6800025}
\textbf{Source problem code:} catalogue entry
\subsection{Definitions and mathematical statement}
For a Riemannian target \((N^n,h)\), a standard closed branched minimal surface is a nonconstant map \(F:(\Sigma,j)\to N\) from a closed connected Riemann surface that is conformal and harmonic, and is an immersion of rank two away from finitely many branch points. Equivalently away from those points its induced surface is minimal; self-intersections and branching are permitted. This makes the usual closed-surface interpretation explicit, but the source does not specify the domain category.

Rivière's question asks for infinitely many distinct such two-dimensional immersions in the intended target \(N^n\). A substantive geometric-distinctness version would require maps \(F_j:\Sigma_j\to N\), \(j\ge1\), with
\[
F_j\text{ branched minimal},\qquad
F_i(\Sigma_i)\ne F_j(\Sigma_j)\quad(i\ne j).
\]
Neither the genus nor a common domain is prescribed. However, distinct maps, maps modulo reparametrization, and distinct images are different requirements: repeated or branched coverings can produce new parametrized maps without new surfaces. The displayed image-distinct version is a candidate meaning, not a silently strengthened source definition.

The record also does not say that \(N\) is closed, complete, or compact, or give a dimension range. Closed-domain existence in arbitrary targets and existence of local minimal disks are different questions. Under image-distinctness, low dimensions require separate attention rather than interpreting the statement as all \(n\) without qualification. These missing target/domain hypotheses prevent a fully quantified universal existence assertion.

The supplied research summary reports the general-codimension question unverified/open and describes active immersed-surface work. It mentions Marques--Neves hypersurface progress, but a hypersurface has dimension \(n-1\); outside ambient dimension three that is not the two-dimensional target here. No general-codimension resolution is inferred from that comparison.
\subsection{Short English statement}
Does the intended Riemannian target contain infinitely many genuinely distinct branched minimal surfaces? The target assumptions, closed-domain convention, dimension range, and treatment of reparametrizations or coverings must be fixed.
\subsection{Sources}
\begin{itemize}
\item Local UnsolvedMath registry, permanent ID \(6800025\), exact source code \texttt{catalogue entry}, ``Minimax minimal surfaces --- Question 2.'' Proposer Tristan Rivière, ETH Zürich, 2018.
\item Morgan and Pansu, \emph{A List of Open Problems in Differential Geometry} (2018), Question environment 25, ``Minimax minimal surfaces''; \url{https://www.imo.universite-paris-saclay.fr/~pansu/problems_MTDG.tex}. Source-TeX extraction; recorded access 2026-07-29.
\item Supplied OPEN-TRIAGE research summary mentions Rivière's programme and Marques--Neves minimal-hypersurface results without complete citations or exact hypotheses.
\end{itemize}
\end{document}
